\documentclass{amsart}
% For dealing with references we use the comment environment
\usepackage{verbatim}
\newenvironment{reference}{\comment}{\endcomment}
%\newenvironment{reference}{}{}
\newenvironment{slogan}{\comment}{\endcomment}
\newenvironment{history}{\comment}{\endcomment}

% For commutative diagrams we use Xy-pic
\usepackage[all]{xy}

% We use 2cell for 2-commutative diagrams.
\xyoption{2cell}
\UseAllTwocells

% We use multicol for the list of chapters between chapters
\usepackage{multicol}

% This is generally recommended for better output
\usepackage{lmodern}
\usepackage[T1]{fontenc}

% For cross-file-references

% Package for hypertext links:
\usepackage{hyperref}

% For any local file, say "hello.tex" you want to link to please

% Theorem environments.
%
\theoremstyle{plain}
\newtheorem{theorem}[subsection]{Theorem}
\newtheorem{proposition}[subsection]{Proposition}
\newtheorem{lemma}[subsection]{Lemma}

\theoremstyle{definition}
\newtheorem{definition}[subsection]{Definition}
\newtheorem{example}[subsection]{Example}
\newtheorem{exercise}[subsection]{Exercise}
\newtheorem{situation}[subsection]{Situation}

\theoremstyle{remark}
\newtheorem{remark}[subsection]{Remark}
\newtheorem{remarks}[subsection]{Remarks}

\numberwithin{equation}{subsection}

% Macros
%
\def\lim{\mathop{\mathrm{lim}}\nolimits}
\def\colim{\mathop{\mathrm{colim}}\nolimits}
\def\Spec{\mathop{\mathrm{Spec}}}
\def\Hom{\mathop{\mathrm{Hom}}\nolimits}
\def\Ext{\mathop{\mathrm{Ext}}\nolimits}
\def\SheafHom{\mathop{\mathcal{H}\!\mathit{om}}\nolimits}
\def\SheafExt{\mathop{\mathcal{E}\!\mathit{xt}}\nolimits}
\def\Sch{\mathit{Sch}}
\def\Mor{\mathop{\mathrm{Mor}}\nolimits}
\def\Ob{\mathop{\mathrm{Ob}}\nolimits}
\def\Sh{\mathop{\mathit{Sh}}\nolimits}
\def\NL{\mathop{N\!L}\nolimits}
\def\CH{\mathop{\mathrm{CH}}\nolimits}
\def\proetale{{pro\text{-}\acute{e}tale}}
\def\etale{{\acute{e}tale}}
\def\QCoh{\mathit{QCoh}}
\def\Ker{\mathop{\mathrm{Ker}}}
\def\Im{\mathop{\mathrm{Im}}}
\def\Coker{\mathop{\mathrm{Coker}}}
\def\Coim{\mathop{\mathrm{Coim}}}

% Boxtimes
%
\DeclareMathSymbol{\boxtimes}{\mathbin}{AMSa}{"02}

%
% Macros for moduli stacks/spaces
%
\def\QCohstack{\mathcal{QC}\!\mathit{oh}}
\def\Cohstack{\mathcal{C}\!\mathit{oh}}
\def\Spacesstack{\mathcal{S}\!\mathit{paces}}
\def\Quotfunctor{\mathrm{Quot}}
\def\Hilbfunctor{\mathrm{Hilb}}
\def\Curvesstack{\mathcal{C}\!\mathit{urves}}
\def\Polarizedstack{\mathcal{P}\!\mathit{olarized}}
\def\Complexesstack{\mathcal{C}\!\mathit{omplexes}}
% \Pic is the operator that assigns to X its picard group, usage \Pic(X)
% \Picardstack_{X/B} denotes the Picard stack of X over B
% \Picardfunctor_{X/B} denotes the Picard functor of X over B
\def\Pic{\mathop{\mathrm{Pic}}\nolimits}
\def\Picardstack{\mathcal{P}\!\mathit{ic}}
\def\Picardfunctor{\mathrm{Pic}}
\def\Deformationcategory{\mathcal{D}\!\mathit{ef}}

\setlength{\emergencystretch}{3em}
\begin{document}
\title[Stacks Project, Tag 0G4D]{474 --- A generic linear combination preserves geometric irreducibility under the stated algebraic-independence hypotheses}
\maketitle
\noindent Copyright (C) 2005--2025 Johan de Jong. Stacks Project authors. Source: \href{https://stacks.math.columbia.edu/tag/0G4D}{Tag 0G4D}.

\noindent Editorial difficulty note (not author text). Difficulty H2: The proof must exclude all new separable algebraic elements over the generic linear-combination field. Translation automorphisms and finiteness of intermediate separable fields force the required invariance, while an additional transcendental parameter handles finite ground fields..

\noindent Editorial provenance note (not author text). History: excerpt from revision \texttt{a0444\allowbreak 6e57e\allowbreak c1fbc\allowbreak 252a8\allowbreak 71afc\allowbreak ec775\allowbreak 2fb28\allowbreak 07b14}, more-morphisms.tex, lines 8906--9009. Reference presentation was adapted; original-author context and core support blocks were retained or added. The mathematical wording is unchanged. Licensed under GNU FDL 1.2 or later; no invariant sections or cover texts. See ../licenses/COPYING. Context blocks reproduce author definitions and supporting results; they are not additional counted problems.

\begin{lemma}
\label{lemma-amazing-bertini-lemma}
\begin{reference}
See pages 71 and 72 of \href{https://stacks.math.columbia.edu/bibliography}{Bibliography: Jou}
\end{reference}
Let $K/k$ be a geometrically irreducible and finitely generated
field extension. Let $n \geq 1$.
Let $g_1, \ldots, g_n \in K$ be elements such that there
exist $c_1, \ldots, c_n \in k$ such that the elements
$$
x_1, \ldots, x_n, \sum g_ix_i, \sum c_ig_i \in K(x_1, \ldots, x_n)
$$
are algebraically independent over $k$. Then
$K(x_1, \ldots, x_n)$ is geometrically irreducible over
$k(x_1, \ldots, x_n, \sum g_ix_i)$.
\end{lemma}

\begin{proof}
Let $c_1, \ldots, c_n \in k$ be as in the statement of the lemma.
Write $\xi = \sum g_ix_i$ and $\delta = \sum c_ig_i$.
For $a \in k$ consider the automorphism $\sigma_a$
of $K(x_1, \ldots, x_n)$ given by the identity on $K$ and the rules
$$
\sigma_a(x_i) = x_i + a c_i
$$
Observe that $\sigma_a(\xi) = \xi + a \delta$ and $\sigma_a(\delta) = \delta$.
Consider the tower of fields
$$
K_0 = k(x_1, \ldots, x_n) \subset
K_1 = K_0(\xi) \subset
K_2 = K_0(\xi, \delta) \subset K(x_1, \ldots, x_n) = \Omega
$$
Observe that $\sigma_a(K_0) = K_0$ and $\sigma_a(K_2) = K_2$.
Let $\theta \in \Omega$ be separable algebraic over $K_1$.
We have to show $\theta \in K_1$, see Algebra, Lemma
\href{https://stacks.math.columbia.edu/tag/0G33}{lemma geometrically irreducible separable elements (Tag 0G33)}.

\medskip\noindent
Denote $K'_2$ the separable algebraic closure of $K_2$ in $\Omega$.
Since $K'_2/K_2$ is finite (Algebra, Lemma
\href{https://stacks.math.columbia.edu/tag/037Q}{lemma make geometrically irreducible (Tag 037Q)}) and
separable there are only a finite number of fields in
between $K'_2$ and $K_2$ (Fields, Lemma \href{https://stacks.math.columbia.edu/tag/030N}{lemma primitive element (Tag 030N)}).
If $k$ is infinite\footnote{We will deal with the finite field case in the
last paragraph of the proof.}, then we can find distinct elements
$a_1, a_2$ of $k$ such that
$$
K_2(\sigma_{a_1}(\theta)) = K_2(\sigma_{a_2}(\theta))
$$
as subfields of $\Omega$. Write $\theta_i = \sigma_{a_i}(\theta)$
and $\xi_i = \sigma_{a_i}(\xi) = \xi + a_i \delta$. Observe that
$$
K_2 = K_0(\xi_1, \xi_2)
$$
as we have $\xi_i = \xi + a_i \delta$,
$\xi = (a_2 \xi_1 - a_1 \xi_2)/(a_2 - a_1)$, and
$\delta = (\xi_1 - \xi_2)/(a_1 - a_2)$.
Since $K_2/K_0$ is purely transcendental of degree $2$ we conclude
that $\xi_1$ and $\xi_2$ are algebraically independent over $K_0$.
Since $\theta_1$ is algebraic over $K_0(\xi_1)$ we conclude that
$\xi_2$ is transcendental over $K_0(\xi_1, \theta_1)$.

\medskip\noindent
By assumption $K/k$ is geometrically irreducible. This implies
that $K(x_1, \ldots, x_n)/K_0$ is geometrically irreducible
(Algebra, Lemma
\href{https://stacks.math.columbia.edu/tag/0G31}{lemma geometrically irreducible base change transcendental (Tag 0G31)}).
This in turn implies that $K_0(\xi_1, \theta_1)/K_0$
is geometrically irreducible as a subextension
(Algebra, Lemma \href{https://stacks.math.columbia.edu/tag/037N}{lemma subalgebra geometrically irreducible (Tag 037N)}).
Since $\xi_2$ is transcendental over $K_0(\xi_1, \theta_1)$
we conclude that $K_0(\xi_1, \xi_2, \theta_1)/K_0(\xi_2)$
is geometrically irreducible (Algebra, Lemma
\href{https://stacks.math.columbia.edu/tag/0G32}{lemma geometrically irreducible add transcendental (Tag 0G32)}).
By our choice of $a_1, a_2$ above we have
$$
K_0(\xi_1, \xi_2, \theta_1) =
K_2(\sigma_{a_1}(\theta)) =
K_2(\sigma_{a_2}(\theta)) =
K_0(\xi_1, \xi_2, \theta_2)
$$
Since $\theta_2$ is separably algebraic over $K_0(\xi_2)$
we conclude by Algebra, Lemma
\href{https://stacks.math.columbia.edu/tag/0G33}{lemma geometrically irreducible separable elements (Tag 0G33)} again that
$\theta_2 \in K_0(\xi_2)$. Taking $\sigma_{a_2}^{-1}$
of this relation givens $\theta \in K_0(\xi) = K_1$ as desired.

\medskip\noindent
This finishes the proof in case $k$ is infinite. If $k$ is finite,
then we can choose a variable $t$ and consider the extension
$K(t)/k(t)$ which is geometrically irreducible by
Algebra, Lemma
\href{https://stacks.math.columbia.edu/tag/0G31}{lemma geometrically irreducible base change transcendental (Tag 0G31)}.
Since it is still be true that
$x_1, \ldots, x_n, \sum g_ix_i, \sum c_ig_i$
in $K(t, x_1, \ldots, x_n)$ are algebraically independent over $k(t)$
we conclude that $K(t, x_1, \ldots, x_n)$
is geometrically irreducible over
$k(t, x_1, \ldots, x_n, \sum g_ix_i)$
by the argument already given.
Then using Algebra, Lemma
\href{https://stacks.math.columbia.edu/tag/0G31}{lemma geometrically irreducible base change transcendental (Tag 0G31)}
once more finishes the job.
\end{proof}
\end{document}
